\documentclass[../../main.tex]{subfiles}
\begin{document}

{\bf [解]}
与えられた関数
\begin{align}
 f(x)=1-\cos x-x\sin x \label{eq:1}
\end{align}
について考える．

\paragraph{(1)}

一階微分を計算すると
\begin{align}
 f'(x)=\sin x -\sin x -x\cos x = -x\cos x \label{eq:2}
\end{align}
である．よって$f(x)$の増減表は以下のようになる．

\begin{table}[htb]
\centering
\caption{$f(x)$ の増減表（$0 \le x \le \pi$）}
\label{tab:f_x_trig}
\begin{tabular}{c|ccccccc}
$x$ & $0$ & $\cdots$ & $\dfrac{\pi}{2}$ & $\cdots$ & $\alpha$ & $\cdots$ & $\pi$ \\ \hline
$f'(x)$ & $0$ & $-$ & $0$ & $+$ & $+$ & $+$ & $\pi$ \\ \hline
$f(x)$ & $0$ & $\searrow$ & $1 - \dfrac{\pi}{2}$（極小） & $\nearrow$ & $0$ & $\nearrow$ & $2$
\end{tabular}
\end{table}

$f(0)=0$かつ$(0,\pi/2)$で$f$は単調減少だから$(0,\pi/2)$に解はなく，$[\pi/2,\pi]$では$f$が$1-\pi/2<0$から$2>0$へ単調増加するので中間値の定理よりただ1つの解$\alpha\in(\pi/2,\pi)$を持つ．よって$0<x<\pi$で$f(x)=0$はただ1つの実解を持つ．以上で題意は示された．

\begin{figure}[htb]
\begin{center}
\begin{tikzpicture}[scale=2.0, >=stealth]

  % 定数設定 (pi ≒ 3.14159, alpha ≒ 2.028757 rad ≒ 116.24 deg)
  \def\piVal{3.14159}
  \def\halfPi{1.570796}
  \def\alphaVal{2.331122}
  \def\threeQuarterPi{2.356194}

  % 座標点
  \coordinate (O) at (0, 0);
  \coordinate (MinPt) at (\halfPi, {1 - \halfPi});
  \coordinate (AlphaPt) at (\alphaVal, 0);
  \coordinate (PiPt) at (\piVal, 2);

  % --- 1. 積分領域のハッチング (J の内訳) ---
  % [0, alpha]: f(x) <= 0
  \fill[pattern=north east lines, pattern color=blue!60]
    plot[domain=0:\alphaVal, samples=60] (\x, {1 - cos(\x r) - \x*sin(\x r)}) -- cycle;

  % [alpha, pi]: f(x) >= 0
  \fill[pattern=north west lines, pattern color=red!60]
    (\alphaVal, 0) -- plot[domain=\alphaVal:\piVal, samples=60] (\x, {1 - cos(\x r) - \x*sin(\x r)}) 
    -- (\piVal, 0) -- cycle;

  % --- 2. 座標軸 ---
  \draw[->, thick] (-0.3, 0) -- (3.6, 0) node[right] {$x$};
  \draw[->, thick] (0, -0.9) -- (0, 2.4) node[above] {$y$};
  \node[below left, font=\small] at (O) {$O$};

  % --- 3. グラフ y = f(x) の描画 ---
  \draw[ultra thick, blue!80!black, domain=0:\piVal, samples=100]
    plot (\x, {1 - cos(\x r) - \x*sin(\x r)})
    node[above left, font=\small] {$y = f(x)$};

  % --- 4. ガイド線と各点のプロット ---
  % 極小点 x = pi/2
  \draw[dashed, gray] (\halfPi, 0) node[above, black, font=\small] {$\dfrac{\pi}{2}$} -- (MinPt);
  \draw[dashed, gray] (0, {1 - \halfPi}) node[left, black, font=\small] {$1-\dfrac{\pi}{2}$} -- (MinPt);
  \fill[blue!80!black] (MinPt) circle (1.2pt) node[below right, font=\scriptsize] {極小};

  % 解 x = alpha
  \fill[purple!80!black] (AlphaPt) circle (1.5pt) 
    node[above left, font=\small, text=purple!80!black] {$\alpha$};

  % 評価点 x = 3pi/4 (参考点)
  \draw[dashed, gray!60] (\threeQuarterPi, 0) node[below, font=\scriptsize, gray!80!black] {$\dfrac{3\pi}{4}$}
    -- (\threeQuarterPi, {1 - cos(\threeQuarterPi r) - \threeQuarterPi*sin(\threeQuarterPi r)});
  \fill[gray!80!black] (\threeQuarterPi, {1 - cos(\threeQuarterPi r) - \threeQuarterPi*sin(\threeQuarterPi r)}) circle (1.0pt);

  % 端点 x = pi
  \draw[dashed, gray] (\piVal, 0) node[below, black, font=\small] {$\pi$} -- (PiPt);
  \draw[dashed, gray] (0, 2) node[left, black, font=\small] {$2$} -- (PiPt);
  \fill[black] (PiPt) circle (1.2pt);

\end{tikzpicture}
\end{center}
\caption{$f(x) = 1 - \cos x - x\sin x$（$0 \le x \le \pi$）の概形と解 $\alpha$}
\end{figure}

\paragraph{(2)}

(1)より$0\le x\le\alpha$で$f(x)\le0$，$\alpha\le x\le\pi$で$f(x)\ge0$だから$J(x)$は
\begin{align}
 J=-\int_0^\alpha f(x)dx+\int_\alpha^\pi f(x)dx \label{eq:3}
\end{align}
と書ける．
$f$の原始関数の1つを$F(x)$とする：
\begin{align}
F(x)
  &=\int f(x)dx \\
  &=x-\sin x-(-x\cos x+\sin x) \\
  &=x(1+\cos x)-2\sin x \label{eq:4}
\end{align}
であり，$F(0)=0$，$F(\pi)=\pi(1+(-1))-0=0$だから\cref{eq:3}に代入して
\begin{align}
 J
  &=F(\pi)+F(0)-2F(\alpha) \\
  &=-2F(\alpha) \\
  &=-2\alpha(1+\cos\alpha)+4\sin\alpha \label{eq:5}
\end{align}
である．
さて，$\alpha$は$f(\alpha)=0$を満たすから\cref{eq:1}より
\begin{align}
 \cos\alpha=1-\alpha\sin\alpha \label{eq:6}
\end{align}
である．\cref{eq:6}を\cref{eq:5}に代入して
\begin{align}
 J
  &=-2\alpha(2-\alpha\sin\alpha)+4\sin\alpha \\
  &=2\alpha^2\sin\alpha-4\alpha+4\sin\alpha \label{eq:7}
\end{align}
である．

\paragraph{(3)}

新しく
\begin{align}
 K=\dfrac{J-\sqrt{2}}{\sqrt{2}} \label{eq:8}
\end{align}
とおいて，この正負を考えれば良い．
\cref{eq:7}を代入して
\begin{align}
 K =\sqrt{2}\alpha^2\sin\alpha-2\sqrt{2}\alpha+2\sqrt{2}\sin\alpha-1 \label{eq:9}
\end{align}
さて，$\sin^2\alpha+\cos^2\alpha=1$に\cref{eq:6}を代入すると
\begin{align}
 \sin^2\alpha+(1-\alpha\sin\alpha)^2=1 \\
 (\alpha^2+1)\sin^2\alpha-2\alpha\sin\alpha=0
\end{align}
である．
$0<\alpha<\pi$より$\sin\alpha\neq0$なので，両辺を$\sin\alpha$で割って
\begin{align}
(\alpha^2+1)\sin\alpha=2\alpha \label{eq:10}
\end{align}
を得る．

\cref{eq:10}を\cref{eq:9}に代入すると
\begin{align}
 K
  &=\sqrt{2}\bigl[(\alpha^2+1)\sin\alpha-2\alpha\bigr]+\sqrt{2}\sin\alpha-1 \\
  &=\sqrt{2}\sin\alpha-1 \label{eq:11}
\end{align}
である．

さて，(1)の増減表と
\begin{align}
 f\Bigl(\dfrac{3\pi}{4}\Bigr)=1+\dfrac{\sqrt{2}}{2}-\dfrac{3\sqrt{2}}{8}\pi>0
\end{align}
より，$\alpha$は
\begin{align}
 \dfrac{\pi}{2} < \alpha < \dfrac{3\pi}{4}
\end{align}
を満たすから，\cref{eq:11}は下から評価できて
\begin{align}
 K > \sqrt{2}\sin\dfrac{3\pi}{4} - 1  = 0 
\end{align}
だから\cref{eq:8}より
\begin{align}
 J > \sqrt{2}
\end{align}
である．

\newpage

\end{document}
